% ====================================================================
% APPENDIX A: STRUCTURAL FRAGILITY OF LOCAL OPTIMIZATION
% ====================================================================
\section{Appendix A: Structural Fragility of Local Optimization}
\label{app:fragility}

Proposition~\ref{thm:nogo} establishes that local optimization requires cells to
maintain precise coupling between metabolic and wave costs that depends on
global network parameters. Here we formalize why such local coupling is
structurally fragile and evolutionarily disfavored.

\subsection{The Sensitivity Amplification Problem}

Let the local coupling parameter $a$ relate to the emergent branching exponent
$\alpha^*$ via the vessel-wall metabolic optimality condition:
\begin{equation}
    a = \frac{1}{1-(\alpha^*)^2} - 2 \quad \implies \quad \alpha^* = \sqrt{1 - \frac{1}{a+2}}.
\end{equation}
We treat this relation as a phenomenological coupling \emph{ansatz} rather
than a first-principles result: it is the simplest single-parameter map with the
required fixed point and a pole near the physiological exponent, and we use it
only to expose the sensitivity-amplification mechanism. The conclusion below
(extreme cumulative sensitivity of local control) is qualitatively robust to the
precise functional form, since any smooth map with a near-pole in this range
produces a large $\sigma$. Differentiating with respect to $a$ yields the
sensitivity coefficient:
\begin{equation}
    \sigma \equiv \frac{1}{\alpha^*}\frac{\partial \alpha^*}{\partial a} = \frac{1}{2(a+2)(a+1)}.
\end{equation}
For a porcine coronary tree where $\alpha^* \approx \VarAlphaStar$, the matching
parameter is $a \approx \VarSensitivityA$, yielding $\sigma \approx
\VarSensitivitySigma$.

\textbf{Cumulative error amplification.} Over a vascular tree of depth $G =
\VarG$, small perturbations in local sensing accumulate multiplicatively over
the $G-1$ internal junctions:
\begin{equation}
    \frac{\delta \alpha^*}{\alpha^*} \approx (G-1) \cdot \sigma \cdot \delta a \approx 27.7 \cdot \delta a.
\end{equation}
To maintain architectural stability within $\delta\alpha^*/\alpha^* < 0.05$ (5\%
tolerance), the required local precision is $\delta a <
\VarSensitivityDeltaA$---a demanding constraint given physiological noise in
shear stress and pressure over the cardiac cycle.

\subsection{Why the Minimax Avoids This Fragility}

In contrast, the network-level minimax requires no local parameter estimation of
global invariants. The saddle point $(\alpha^*, \eta^*)$ emerges from the
\emph{topological structure} of the optimization landscape: it is the unique
point where metabolic cost equals wave cost, determined entirely by
dimensionless structural parameters $(G, N, p, \alpha_w)$.

Mechanotransduction-driven remodeling based on local shear stress $\tau$ and
circumferential strain $\epsilon$ converges spontaneously to this attractor
because the equal-cost condition is a stationary point of the global energy
landscape. The network does not \emph{compute} the optimal state; it
\emph{relaxes} toward it via gradient descent on a robust cost functional.

\textbf{Robustness vs fragility.} Local optimization requires fine-tuned
sensitivity to global parameters and is vulnerable to cumulative error
amplification. Network-level minimax optimization is topologically robust: the
attractor $\alpha^* \approx \VarAlphaStar$ emerges from structural constraints
and remains stable across seven orders of magnitude in body mass, despite
massive variations in metabolic rate, blood viscosity, and cardiac output.

This structural robustness explains why evolution favors the minimax solution:
it achieves near-optimal performance without requiring cells to "know" the
global scale of the organism.
